% ============================================================
%  Apprendre Go : Pourquoi ? Pour Quoi Faire ?
%  Auteur  : Axel Konan
%  Moteur  : pdflatex -shell-escape formation-go.tex (x2)
% ============================================================
\documentclass[11pt,a4paper,openany]{book}

% --- Encodage ------------------------------------------------
\usepackage[utf8]{inputenc}
\usepackage[T1]{fontenc}
% \usepackage{microtype} % desactive pour compat

% --- Couleurs ------------------------------------------------
\usepackage[dvipsnames,svgnames,x11names]{xcolor}
\definecolor{GoBlue}      {HTML}{00ACD7}
\definecolor{GoDark}      {HTML}{0D1B2A}
\definecolor{GoAccent}    {HTML}{FF6D00}
\definecolor{GoLight}     {HTML}{E8F7FB}
\definecolor{GoGrey}      {HTML}{546E7A}
\definecolor{GoLightGrey} {HTML}{F4F7FB}
\definecolor{CodeBg}      {HTML}{1E2B38}
\definecolor{CodeFg}      {HTML}{CDD6F4}
\definecolor{GreenCode}   {HTML}{2ECC71}
\definecolor{OrangeWarn}  {HTML}{F39C12}

% --- Geometrie -----------------------------------------------
\usepackage[a4paper,top=24mm,bottom=24mm,
            inner=28mm,outer=22mm,
            headsep=8mm,footskip=12mm]{geometry}

% --- TikZ ---------------------------------------------------
\usepackage{tikz}
\usetikzlibrary{shapes,arrows.meta,positioning,calc,decorations.pathreplacing}

% --- En-tetes et pieds de page ------------------------------
\usepackage{fancyhdr}
\pagestyle{fancy}
\fancyhf{}
\fancyhead[LE]{\small\color{GoGrey}\leftmark}
\fancyhead[RO]{\small\color{GoGrey}\rightmark}
\fancyfoot[LE,RO]{\small\color{GoGrey}\thepage}
\fancyfoot[CE,CO]{\small\color{GoBlue}%
  Apprendre Go : Pourquoi~? Pour Quoi Faire~?}
\renewcommand{\headrulewidth}{0.4pt}
\renewcommand{\footrulewidth}{0.4pt}

% --- Titres -------------------------------------------------
\usepackage{titlesec}

\titleformat{\chapter}[display]
  {\normalfont\bfseries\color{GoDark}}
  {\begin{tikzpicture}
     \fill[GoBlue](0,0) rectangle (\textwidth,1cm);
     \node[anchor=east,text=white,font=\large\bfseries]
       at (\textwidth-5mm,0.5cm)
       {\chaptertitlename\ \thechapter};
   \end{tikzpicture}}
  {4pt}
  {\LARGE\bfseries\color{GoDark}}
  [\vspace{1mm}{\color{GoBlue}\hrule height 1.5pt}\vspace{3mm}]

\titleformat{\section}
  {\large\bfseries\color{GoDark}}
  {\color{GoBlue}\thesection}{0.8em}{}
  [\vspace{0.5mm}{\color{GoLightGrey}\hrule height 2pt}\vspace{1mm}]

\titleformat{\subsection}
  {\normalsize\bfseries\color{GoDark}}
  {\color{GoAccent}\thesubsection}{0.8em}{}

\titlespacing*{\chapter}   {0pt}{0pt}{14pt}
\titlespacing*{\section}   {0pt}{12pt}{5pt}
\titlespacing*{\subsection}{0pt}{8pt}{3pt}

% --- Table des matieres -------------------------------------
\usepackage[titles]{tocloft}
\renewcommand{\cftchapfont}{\bfseries\color{GoDark}}
\renewcommand{\cftsecfont}{\color{GoGrey}}
\renewcommand{\cftchappagefont}{\bfseries\color{GoBlue}}
\renewcommand{\cftsecpagefont}{\color{GoGrey}}
\setlength{\cftbeforechapskip}{4pt}

% --- Boites tcolorbox ----------------------------------------
\usepackage{tcolorbox}
\tcbuselibrary{skins,breakable}

\newtcolorbox{boiteConseil}{
  enhanced,breakable,
  colback=GoLight,colframe=GoBlue,
  leftrule=4pt,toprule=0pt,bottomrule=0pt,rightrule=0pt,
  arc=2pt,left=6pt,right=6pt,top=4pt,bottom=4pt,
  title={\small\bfseries Astuce},
  fonttitle=\color{white},
  attach boxed title to top left={yshift=-2.5mm,xshift=4mm},
  boxed title style={colback=GoBlue,colframe=GoBlue,arc=2pt}}

\newtcolorbox{boiteAttention}{
  enhanced,breakable,
  colback={orange!8},colframe=OrangeWarn,
  leftrule=4pt,toprule=0pt,bottomrule=0pt,rightrule=0pt,
  arc=2pt,left=6pt,right=6pt,top=4pt,bottom=4pt,
  title={\small\bfseries Attention},
  fonttitle=\color{white},
  attach boxed title to top left={yshift=-2.5mm,xshift=4mm},
  boxed title style={colback=OrangeWarn,colframe=OrangeWarn,arc=2pt}}

\newtcolorbox{boiteZoom}[1]{
  enhanced,breakable,
  colback={GoGrey!8},colframe=GoGrey,
  leftrule=4pt,toprule=0pt,bottomrule=0pt,rightrule=0pt,
  arc=2pt,left=6pt,right=6pt,top=4pt,bottom=4pt,
  title={\small\bfseries Zoom profil~--- #1},
  fonttitle=\color{white},
  attach boxed title to top left={yshift=-2.5mm,xshift=4mm},
  boxed title style={colback=GoGrey,colframe=GoGrey,arc=2pt}}

\newtcolorbox{boiteProjet}[1]{
  enhanced,breakable,
  colback={GoAccent!7},colframe=GoAccent,
  arc=4pt,left=8pt,right=8pt,top=5pt,bottom=5pt,
  title={\small\bfseries Projet fil rouge~--- #1},
  fonttitle=\color{white},
  attach boxed title to top left={yshift=-2.5mm,xshift=6mm},
  boxed title style={colback=GoAccent,colframe=GoAccent,arc=3pt}}

\newtcolorbox{boiteCitation}{
  enhanced,
  colback=GoDark,colframe=GoDark,
  arc=4pt,left=10pt,right=10pt,top=7pt,bottom=7pt,
  fontupper=\normalsize\itshape\color{GoLight}}

% --- Code listings -------------------------------------------
\usepackage{listings}

\lstdefinestyle{goStyle}{
  backgroundcolor=\color{CodeBg},
  basicstyle=\small\ttfamily\color{CodeFg},
  keywordstyle=\color{GoBlue}\bfseries,
  stringstyle=\color{GreenCode},
  commentstyle=\color{GoGrey}\itshape,
  numberstyle=\tiny\color{GoGrey},
  numbers=left,numbersep=8pt,
  breaklines=true,frame=none,
  xleftmargin=14pt,xrightmargin=4pt,
  tabsize=4,showstringspaces=false,
  language=[Objective]C,
  morekeywords={func,package,import,var,const,type,struct,
    interface,map,chan,go,defer,select,range,fallthrough,
    nil,true,false,iota,make,new,append,len,cap,close,
    delete,copy,panic,recover,return,if,else,for,switch,
    case,default,break,continue,goto}}

\lstdefinestyle{bashStyle}{
  backgroundcolor=\color{GoDark},
  basicstyle=\small\ttfamily\color{GoLight},
  numbers=none,breaklines=true,frame=none,
  xleftmargin=10pt,showstringspaces=false}

\lstset{style=goStyle,aboveskip=6pt,belowskip=6pt}

\newcommand{\code}[1]{%
  \colorbox{GoLightGrey}{\small\ttfamily\color{GoDark}#1}}

% --- Tableaux ------------------------------------------------
\usepackage{booktabs,tabularx,colortbl}
\newcolumntype{L}{>{\raggedright\arraybackslash}X}
\newcolumntype{C}{>{\centering\arraybackslash}X}
\newcommand{\tableauStyle}{%
  \rowcolors{2}{GoLightGrey}{white}\arrayrulecolor{GoBlue}}

% --- Listes --------------------------------------------------
\usepackage{enumitem}
\setlist[itemize]{leftmargin=1.5em,itemsep=2pt,topsep=3pt,
  label={\color{GoBlue}\textbullet}}
\setlist[enumerate]{leftmargin=1.5em,itemsep=2pt,topsep=3pt,
  label={\color{GoBlue}\arabic*.}}

% --- Liens et PDF --------------------------------------------
\usepackage[colorlinks=true,linkcolor=GoBlue,urlcolor=GoAccent,
  citecolor=GoGrey,
  pdftitle={Apprendre Go : Pourquoi ? Pour Quoi Faire ?},
  pdfauthor={Axel Konan},
  pdfsubject={Formation Go complete -- 5 modules}]{hyperref}

% --- Divers --------------------------------------------------
\usepackage{graphicx,float,caption,setspace}
\captionsetup{font=small,labelfont={bf,color=GoBlue},textfont=it}
\setstretch{1.2}
\setlength{\parindent}{0pt}
\setlength{\parskip}{5pt}

% --- Macros utilitaires --------------------------------------
\newcommand{\moduleHeader}[3]{%
  \vspace{4mm}%
  \begin{tikzpicture}
    \fill[#3](0,0) rectangle (\textwidth,0.9cm);
    \node[anchor=west,text=white,font=\large\bfseries]
      at (5mm,0.45cm) {MODULE #1};
    \node[anchor=east,text=white,font=\normalsize]
      at (\textwidth-5mm,0.45cm) {#2};
  \end{tikzpicture}\vspace{5mm}}

% --- Page de titre -------------------------------------------
\newcommand{\faireTitre}{%
\begin{titlepage}
\begin{tikzpicture}[remember picture,overlay]
  % Fond sombre
  \fill[GoDark]
    (current page.south west) rectangle (current page.north east);
  % Bande haute bleue
  \fill[GoBlue]
    (current page.north west)
    rectangle ([yshift=-1.4cm]current page.north east);
  % Bande basse orange
  \fill[GoAccent]
    (current page.south west)
    rectangle ([yshift=0.9cm]current page.south east);
  % Cercle decoratif droit
  \fill[GoBlue!30]
    ([xshift=13cm,yshift=-5.5cm]current page.north west) circle (4cm);
  \fill[GoDark]
    ([xshift=13cm,yshift=-5.5cm]current page.north west) circle (2.6cm);
  \node[text=GoBlue,font=\fontsize{52}{52}\selectfont\bfseries]
    at ([xshift=13cm,yshift=-5.5cm]current page.north west) {Go};
  % Titre
  \node[anchor=west,text=white,
    font=\fontsize{36}{42}\selectfont\bfseries,text width=9cm]
    at ([xshift=1.5cm,yshift=-3cm]current page.north west)
    {Apprendre Go :};
  % Sous-titre
  \node[anchor=west,text=GoBlue,
    font=\fontsize{17}{22}\selectfont\bfseries,text width=9cm]
    at ([xshift=1.5cm,yshift=-5.1cm]current page.north west)
    {Mais Pourquoi ?\\[2mm]Et Pour Quoi Faire ?};
  % Description
  \node[anchor=west,text=GoGrey!80,
    font=\fontsize{11}{15}\selectfont\itshape,text width=9cm]
    at ([xshift=1.5cm,yshift=-7.1cm]current page.north west)
    {De la philosophie du langage\\
     a l'architecture de systemes distribues.\\[1mm]
     5 modules $\cdot$ 13 chapitres $\cdot$ 2 projets fil rouge.};
  % Separateur
  \draw[GoBlue,line width=1.2pt]
    ([xshift=1.5cm,yshift=-8.5cm]current page.north west) --
    ([xshift=10.5cm,yshift=-8.5cm]current page.north west);
  % Auteur
  \node[anchor=west,text=white,
    font=\fontsize{15}{15}\selectfont\bfseries]
    at ([xshift=1.5cm,yshift=-9.4cm]current page.north west)
    {Axel Konan};
  \node[anchor=west,text=GoGrey!60,
    font=\fontsize{10}{10}\selectfont]
    at ([xshift=1.5cm,yshift=-10.3cm]current page.north west)
    {Architecte Logiciel $\cdot$ Formateur $\cdot$ Vulgarisateur Technique};
  % Badges modules
  \foreach \xpos/\lbl/\col in {
    1.5/Module 01/GoBlue,
    4.2/Module 02/GoBlue!80,
    6.9/Module 03/GoAccent,
    9.6/Module 04/GoAccent!70,
    12.3/Module 05/GoGrey}{
    \fill[\col]
      ([xshift=\xpos cm,yshift=4.8cm]current page.south west)
      rectangle ([xshift=\xpos cm+2.4cm,yshift=5.5cm]current page.south west);
    \node[text=white,font=\fontsize{7}{7}\selectfont\bfseries]
      at ([xshift=\xpos cm+1.2cm,yshift=5.15cm]current page.south west) {\lbl};}
  % Version
  \node[anchor=east,text=GoGrey!40,font=\fontsize{8}{8}\selectfont]
    at ([xshift=-1cm,yshift=1.4cm]current page.south east)
    {v1.0 $\cdot$ \the\year{} $\cdot$ CC BY-NC-SA 4.0};
\end{tikzpicture}
\end{titlepage}}

% ============================================================
%  DEBUT DU DOCUMENT
% ============================================================
\begin{document}
\faireTitre

% --- Liminaires ----------------------------------------------
\frontmatter
\pagestyle{plain}

\clearpage
\vspace*{\fill}
\begin{center}
  {\Large\bfseries\color{GoDark}%
    Apprendre Go : Mais Pourquoi ? Et Pour Quoi Faire ?}\\[5mm]
  {\color{GoGrey}Formation complete~---
    5 modules $\cdot$ 13 chapitres $\cdot$ 2 projets fil rouge}\\[8mm]
  \begin{tabular}{rl}
    \textbf{Auteur}  & Axel Konan \\
    \textbf{Version} & 1.0 \\
    \textbf{Annee}   & \the\year \\
    \textbf{Depot GitHub} &
      \href{https://github.com/votre-pseudo/formation-go}%
           {github.com/votre-pseudo/formation-go} \\
    \textbf{Licence} & Creative Commons CC BY-NC-SA 4.0 \\
  \end{tabular}\\[8mm]
  {\color{GoBlue}\hrule}\vspace{4mm}
  {\small\color{GoGrey}
    Ce document est distribue gratuitement.\\
    Partagez-le et adaptez-le en citant l'auteur.}
\end{center}
\vspace*{\fill}

\tableofcontents
\clearpage

% --- PREFACE -------------------------------------------------
\chapter*{Preface}
\addcontentsline{toc}{chapter}{Preface}

\begin{boiteCitation}
``La simplicite est la sophistication supreme.''\\
\hfill --- Rob Pike, co-createur de Go
\end{boiteCitation}

\vspace{4mm}

Il existe des dizaines de ressources pour apprendre Go.
Alors pourquoi celle-ci ?

Parce que toutes ces ressources repondent a la question \textbf{Comment~?}
et presque aucune ne repond d'abord a la question \textbf{Pourquoi~?}

Sans le Pourquoi, le Comment ne tient pas. On apprend des syntaxes
qu'on oublie. On abandonne au troisieme chapitre parce qu'on ne voit
pas ou ca mene.

\begin{boiteCitation}
Moins de magie. Plus de clarte.\\
Des systemes qui tiennent debout.
\end{boiteCitation}

\medskip
Chaque chapitre suit toujours le meme schema~:
\begin{enumerate}
  \item \textbf{Le probleme}~--- La frustration concrete que ce chapitre resout
  \item \textbf{L'intuition}~--- L'analogie ou l'image mentale avant le code
  \item \textbf{La solution Go}~--- Le code commente, explique ligne par ligne
  \item \textbf{Ce qu'il faut retenir}~--- L'essentiel en 3 points maximum
  \item \textbf{Pour aller plus loin}~--- Des ressources optionnelles
\end{enumerate}

% --- INTRODUCTION --------------------------------------------
\chapter*{Introduction}
\addcontentsline{toc}{chapter}{Introduction}

\begin{boiteProjet}{gowatch --- Outil CLI de monitoring systeme}
Un binaire unique, concurrent, sans dependances externes.
Du simple affichage du Module~01 jusqu'au scanner reseau securise
du Module~05. Un seul executable, zero dependance.
\end{boiteProjet}

\begin{boiteProjet}{gohub --- Dashboard API REST}
Une API HTTP testee, dockerisee, en image Docker de moins de 10~Mo.
Introduite au Module~04, elle s'appuie sur tout ce qu'on a construit avant.
\end{boiteProjet}

\section*{Parcours recommandes par profil}

\begin{center}
\tableauStyle
\begin{tabularx}{\textwidth}{>{\bfseries}l L l}
\toprule
\rowcolor{GoDark}
\color{white}Profil &
\color{white}Objectif &
\color{white}Parcours \\
\midrule
Backend / Fullstack & API performante en production    & 01 a 04 \\
DevOps / Cloud      & Outils CLI + Cloud Native        & 01 a 05 \\
Python / JS         & Memoire + concurrence            & 01 a 04 \\
Low-Level / Secu    & Scanner reseau + cryptographie   & 01, 02, 03, 05 \\
\bottomrule
\end{tabularx}
\end{center}

% ============================================================
%  CORPS
% ============================================================
\mainmatter
\pagestyle{fancy}

% ============================================================
\moduleHeader{01}{L'Eveil --- Pourquoi Go existe}{GoBlue}
% ============================================================

\chapter{Le syndrome de la complexite}

\begin{boiteCitation}
``Complexity is the enemy of reliability.''
\hfill --- Rob Pike
\end{boiteCitation}

\section{Le probleme}

Nous sommes en 2007, dans les bureaux de Google. Trois langages dominent~:

\begin{itemize}
  \item \textbf{C++}~: rapide, mais compile en 45 minutes. Gestion memoire dangereuse.
  \item \textbf{Java}~: accessible, mais verbeux, lent au demarrage, gourmand en RAM.
  \item \textbf{Python}~: elegant, mais s'effondre sous la charge concurrente.
\end{itemize}

\textbf{Robert Griesemer, Rob Pike et Ken Thompson} decident de creer Go.
Pas pour le plaisir de l'exercice~--- pour resoudre des problemes reels
a grande echelle.

\section{Les 3 piliers fondateurs}

\subsection{Pilier 1 --- La Simplicite}

Pas de classes. Pas d'heritage. Pas de surcharge d'operateurs.
Ce n'est pas un manque~--- c'est une decision deliberee.

\begin{boiteConseil}
\code{gofmt} formate automatiquement votre code selon les conventions
officielles. Pas de debat tabs vs espaces. Pas de style d'equipe a negocier.
\code{gofmt} a tranche une fois pour toutes.
\end{boiteConseil}

\subsection{Pilier 2 --- L'Efficacite}

Un \textbf{binaire natif unique}, sans dependance externe. Demarre en
millisecondes. Une fraction de la memoire d'une application Java equivalente.
La compilation d'un projet moyen prend quelques \textbf{secondes}, pas des minutes.

\subsection{Pilier 3 --- La Concurrence Native}

\begin{boiteCitation}
``Don't communicate by sharing memory;\\
share memory by communicating.''
\hfill --- Rob Pike
\end{boiteCitation}

\section{Comparaison honnete des langages}

\begin{center}
\tableauStyle
\begin{tabularx}{\textwidth}{L C C C C C}
\toprule
\rowcolor{GoDark}
\color{white}Critere &
\color{white}Go &
\color{white}Python &
\color{white}Java &
\color{white}C++ &
\color{white}Rust \\
\midrule
Vitesse execution      & 4/5 & 2/5 & 3/5 & 5/5 & 5/5 \\
Vitesse compilation    & 5/5 & N/A & 2/5 & 1/5 & 2/5 \\
Concurrence native     & 5/5 & 2/5 & 3/5 & 2/5 & 4/5 \\
Facilite apprentissage & 4/5 & 5/5 & 3/5 & 1/5 & 2/5 \\
Deploiement binaire    & 5/5 & 2/5 & 2/5 & 4/5 & 5/5 \\
\bottomrule
\end{tabularx}
\end{center}

\chapter{La Zen Attitude du Gopher}

\section{Installation en 3 commandes}

\begin{lstlisting}[style=bashStyle]
wget https://go.dev/dl/go1.23.0.linux-amd64.tar.gz
sudo tar -C /usr/local -xzf go1.23.0.linux-amd64.tar.gz
export PATH=$PATH:/usr/local/go/bin
go version   # go version go1.23.0 linux/amd64
\end{lstlisting}

\section{La toolchain integree}

\begin{center}
\tableauStyle
\begin{tabularx}{\textwidth}{>{\ttfamily\small}l L}
\toprule
\rowcolor{GoDark}\color{white}Commande & \color{white}Usage \\
\midrule
go run main.go    & Compiler et executer (developpement) \\
go build -o app . & Compiler en binaire (production) \\
go test ./...     & Lancer tous les tests \\
go fmt ./...      & Formatter le code automatiquement \\
go vet ./...      & Analyser les erreurs logiques \\
go mod tidy       & Nettoyer les dependances \\
\bottomrule
\end{tabularx}
\end{center}

\section{Premier programme --- analyse ligne par ligne}

\begin{lstlisting}[style=goStyle]
package main   // "main" = point d'entree du programme

import "fmt"   // Import inutilise = erreur de compilation !

func main() {  // Lance au demarrage du binaire
    fmt.Println("Bonjour, monde.")
}
\end{lstlisting}

\begin{boiteAttention}
Si vous importez un package et ne l'utilisez pas, \textbf{Go refuse de compiler}.
Frustrant au debut --- libérateur sur le long terme. Pas de code mort accumule.
\end{boiteAttention}

\begin{boiteProjet}{gowatch v0.1 --- Premier binaire autonome}
\begin{lstlisting}[style=goStyle]
package main
import ("fmt"; "runtime")
func main() {
    fmt.Println("=== gowatch v0.1 ===")
    fmt.Printf("Systeme  : %s/%s\n", runtime.GOOS, runtime.GOARCH)
    fmt.Printf("Go       : %s\n",    runtime.Version())
    fmt.Printf("CPU(s)   : %d\n",    runtime.NumCPU())
    fmt.Printf("Goroutines : %d\n",  runtime.NumGoroutine())
}
// Compiler : go build -o gowatch .
// Resultat : un seul fichier executable, zero dependance
\end{lstlisting}
\end{boiteProjet}

% ============================================================
\moduleHeader{02}{La Forge --- Les fondamentaux sans s'ennuyer}{GoBlue!80}
% ============================================================

\chapter{Les briques fondamentales}

\section{Variables et inference de type}

\begin{lstlisting}[style=goStyle]
// Forme longue
var nom  string = "Alice"
var age  int    = 30

// Forme courte -- idiomatique (a l'interieur des fonctions)
nom := "Alice"   // Go infere le type string
age := 30        // Go infere le type int

// Regle d'or : toute variable declaree DOIT etre utilisee
\end{lstlisting}

\section{Slices --- La structure de donnees centrale}

\begin{lstlisting}[style=goStyle]
nombres := []int{10, 20, 30, 40, 50}
nombres  = append(nombres, 60)            // Ajout dynamique

for index, val := range nombres {         // Iteration
    fmt.Printf("%d : %d\n", index, val)
}

donnees := make([]int, 0, 1000)           // Pre-allocation
\end{lstlisting}

\section{Maps --- Le dictionnaire haute performance}

\begin{lstlisting}[style=goStyle]
capitales := map[string]string{
    "France":  "Paris",
    "Senegal": "Dakar",
}

// Toujours verifier l'existence d'une cle !
capitale, existe := capitales["Allemagne"]
if existe {
    fmt.Println("Capitale :", capitale)
} else {
    fmt.Println("Pays non trouve")
}
\end{lstlisting}

\begin{boiteAttention}
Acceder a une cle inexistante retourne la valeur zero du type silencieusement.
Toujours utiliser la forme a deux retours~: \code{val, ok := maMap["cle"]}.
\end{boiteAttention}

\chapter{La logique sans fioritures}

\section{La boucle \texttt{for} universelle}

Go n'a qu'un seul mot-cle de boucle. Il couvre \textit{tous} les cas~:

\begin{lstlisting}[style=goStyle]
for i := 0; i < 5; i++ { }              // Compteur classique
for compteur < 10 { compteur++ }        // Equivalent while
for { if stop { break } }              // Boucle infinie

for index, val := range maSlice { }   // Range slice
for cle, val  := range maMap   { }    // Range map
\end{lstlisting}

\section{Gestion d'erreurs --- Le choix le plus controverse}

En Go, les erreurs sont des \textbf{valeurs}, pas des exceptions~:

\begin{lstlisting}[style=goStyle]
// Une fonction qui peut echouer retourne DEUX valeurs
fichier, err := os.Open("config.json")
if err != nil {
    return fmt.Errorf("impossible d'ouvrir : %w", err)
}
defer fichier.Close()  // Execute a la sortie de la fonction
\end{lstlisting}

\begin{boiteZoom}{Developpeurs Java / Python}
Oui, \code{if err != nil} partout semble verbeux. Mais combien d'exceptions
silencieuses attrapez-vous avec un \code{catch (Exception e) \{\}} vide ?
Go rend l'ignorance d'une erreur visible et deliberee~--- pas accidentelle.
\end{boiteZoom}

\chapter{Composition vs Heritage}

\section{Structs et methodes}

\begin{lstlisting}[style=goStyle]
type Metrique struct {
    Nom    string   // Majuscule = exporte (public)
    Valeur float64
    Unite  string
}

// Methode avec recepteur par valeur (lecture seule)
func (m Metrique) String() string {
    return fmt.Sprintf("%s=%.2f%s", m.Nom, m.Valeur, m.Unite)
}

cpu := Metrique{Nom: "CPU", Valeur: 23.5, Unite: "%"}
fmt.Println(cpu)  // CPU=23.50%  <- String() automatique
\end{lstlisting}

\section{Interfaces --- Duck Typing statique}

\begin{lstlisting}[style=goStyle]
// Definir une interface
type Renderer interface {
    Render(snap Snapshot) error
}

// TextRenderer l'implemente SANS "implements" explicite
type TextRenderer struct{}
func (r TextRenderer) Render(s Snapshot) error { return nil }

// Si votre type a la methode Render() -> il implemente Renderer
// Verifie a la compilation, pas a l'execution
\end{lstlisting}

% ============================================================
\moduleHeader{03}{Le Super-Pouvoir --- La concurrence massive}{GoAccent}
% ============================================================

\chapter{Goroutines --- Le chaos ordonne}

\begin{boiteCitation}
Thread OS = un camion 38 tonnes.\\
Goroutine = un coursier a velo.\\
La ville prefere les velos.
\end{boiteCitation}

\section{Le modele M:N --- Pourquoi c'est si leger}

\begin{center}
\begin{tikzpicture}[
  box/.style={draw=GoBlue,fill=GoLight,rounded corners=3pt,
    minimum width=1.8cm,minimum height=0.65cm,font=\small},
  arr/.style={->,>=Stealth,thick,color=GoBlue}]
\foreach \i/\x in {1/0,2/2.1,3/4.2,4/6.3,5/8.4}{
  \node[box,fill=GoBlue!20] at (\x,3) {\small Go~\i};}
\node[font=\small\color{GoGrey}] at (4.2,2.4)
  {Goroutines multiplexees par le scheduler Go};
\foreach \i/\x in {1/1.5,2/4.7,3/7.8}{
  \node[box,fill=GoAccent!20,draw=GoAccent] at (\x,1.5)
  {\small Thread~\i};}
\node[font=\small\color{GoGrey}] at (4.2,0.9)
  {Threads OS geres par le noyau};
\node[box,fill=GoDark,text=white,minimum width=5cm,draw=GoDark]
  at (4.2,0) {\small CPU~(coeurs physiques)};
\foreach \gx/\tx in {0/1.5,2.1/1.5,4.2/4.7,6.3/7.8,8.4/7.8}{
  \draw[arr](\gx,2.68)--(\tx,1.83);}
\foreach \tx in {1.5,4.7,7.8}{
  \draw[arr](\tx,1.18)--(4.2,0.33);}
\end{tikzpicture}
\end{center}

\begin{itemize}
  \item Goroutine~: \textbf{\raise.17ex\hbox{$\sim$}2~Ko} de memoire
        vs 1--2~Mo pour un thread OS
  \item Creer une goroutine~: quelques \textbf{microsecondes}
  \item On peut en lancer \textbf{des centaines de milliers}
        sur une machine ordinaire
\end{itemize}

\section{Lancer une goroutine}

\begin{lstlisting}[style=goStyle]
func main() {
    go direBonjour("Alice")   // Non bloquant -- arriere-plan
    go direBonjour("Bob")

    // Piege classique : toujours passer i en parametre
    for i := 0; i < 5; i++ {
        go func(n int) { fmt.Println(n) }(i)  // Correct
    }
}
\end{lstlisting}

\begin{boiteAttention}
\textbf{Le piege des closures en boucle}~:
\begin{lstlisting}[style=goStyle]
// MAUVAIS -- toutes les goroutines partagent la meme variable i
for i := 0; i < 5; i++ {
    go func() { fmt.Println(i) }()  // Peut afficher 5,5,5,5,5
}
\end{lstlisting}
\end{boiteAttention}

\chapter{Channels --- L'art de communiquer}

\begin{lstlisting}[style=goStyle]
ch := make(chan int)           // Channel unbuffered

go func() { ch <- 42 }()      // Envoyer -- bloque jusqu'a reception
valeur := <-ch                 // Recevoir -- bloque jusqu'a envoi

// Buffered : n'attend pas immediatement un recepteur
ch2 := make(chan string, 10)

// Select : attendre le premier channel pret
select {
case msg := <-ch1:
    fmt.Println("ch1:", msg)
case msg := <-ch2:
    fmt.Println("ch2:", msg)
case <-time.After(5 * time.Second):
    fmt.Println("timeout !")   // Pattern idiomatique
}
\end{lstlisting}

\begin{boiteProjet}{gowatch v0.6 --- Collecte concurrente avec channels}
\begin{lstlisting}[style=goStyle]
func collecterSources(sources []Source) <-chan Metrique {
    out := make(chan Metrique, len(sources))
    var wg sync.WaitGroup

    for _, src := range sources {
        wg.Add(1)
        s := src
        go func() {
            defer wg.Done()
            val, err := s.Collecte()
            out <- Metrique{Source: s.Nom, Valeur: val}
            _ = err
        }()
    }

    go func() { wg.Wait(); close(out) }()
    return out  // Fan-out / Fan-in en 10 lignes
}
\end{lstlisting}
\end{boiteProjet}

\chapter{Patterns de concurrence avances}

\begin{lstlisting}[style=goStyle]
// WaitGroup : attendre un groupe de goroutines
var wg sync.WaitGroup
wg.Add(1)
go func() { defer wg.Done() /* travail */ }()
wg.Wait()   // Bloque jusqu'a ce que le compteur = 0

// Mutex : proteger un etat partage
var mu sync.Mutex
mu.Lock()
defer mu.Unlock()
compteur++   // Acces exclusif garanti

// Context : propager les annulations en cascade
ctx, cancel := context.WithTimeout(
    context.Background(), 5*time.Second)
defer cancel()   // Toujours appeler cancel() !
\end{lstlisting}

% ============================================================
\moduleHeader{04}{L'Architecte --- Systemes serieux en production}{GoAccent!80}
% ============================================================

\chapter{Cloud Native et Microservices}

\section{Serveur HTTP sans framework}

\begin{lstlisting}[style=goStyle]
mux := http.NewServeMux()
mux.HandleFunc("/health",      handleHealth)
mux.HandleFunc("/api/metrics", handleMetrics)

// TOUJOURS configurer les timeouts en production
srv := &http.Server{
    Addr:         ":8080",
    Handler:      mux,
    ReadTimeout:  5 * time.Second,
    WriteTimeout: 10 * time.Second,
    IdleTimeout:  120 * time.Second,
}
srv.ListenAndServe()
\end{lstlisting}

\begin{boiteConseil}
\code{http.ListenAndServe} n'a pas de timeouts. Une connexion bloquee
peut geler une goroutine indefiniment. Toujours utiliser
\code{http.Server} avec des timeouts explicites.
\end{boiteConseil}

\section{Reponses JSON et middleware}

\begin{lstlisting}[style=goStyle]
// Helper JSON -- zero allocation intermediaire
func jsonResponse(w http.ResponseWriter, code int, v any) {
    w.Header().Set("Content-Type", "application/json")
    w.WriteHeader(code)
    json.NewEncoder(w).Encode(v)  // Ecriture directe
}

// Middleware : Handler -> Handler
func middlewareLogging(next http.Handler) http.Handler {
    return http.HandlerFunc(func(w http.ResponseWriter, r *http.Request) {
        debut := time.Now()
        next.ServeHTTP(w, r)
        slog.Info("http", "method", r.Method,
            "path", r.URL.Path, "dur", time.Since(debut))
    })
}
\end{lstlisting}

\chapter{Tooling et Qualite Production}

\section{Table-Driven Tests --- L'idiome central de Go}

\begin{lstlisting}[style=goStyle]
func TestStore_Historique(t *testing.T) {
    tests := []struct {
        nom     string
        inserts int
        limite  int
        attendu int
    }{
        {"limite < total", 10, 5,  5},
        {"store vide",      0, 5,  0},
        {"limite > total",  3, 10, 3},
    }
    for _, tc := range tests {
        tc := tc   // Capture pour les sous-tests paralleles
        t.Run(tc.nom, func(t *testing.T) {
            s := store.New()
            for i := 0; i < tc.inserts; i++ { s.Ajouter(Snapshot{}) }
            if got := len(s.Historique(tc.limite)); got != tc.attendu {
                t.Errorf("attendu %d, obtenu %d", tc.attendu, got)
            }
        })
    }
}
\end{lstlisting}

\chapter{Deploiement et DevOps}

\section{Docker multi-stage~: de 800 Mo a 8 Mo}

\begin{lstlisting}[style=bashStyle]
# Stage 1 : Compiler dans une image complete
FROM golang:1.23-alpine AS builder
COPY . .
RUN CGO_ENABLED=0 go build -ldflags="-s -w" -o /gohub .

# Stage 2 : Image finale VIDE -- uniquement le binaire
FROM scratch
COPY --from=builder /etc/ssl/certs/ca-certificates.crt /etc/ssl/certs/
COPY --from=builder /gohub /gohub
ENTRYPOINT ["/gohub"]
\end{lstlisting}

\begin{lstlisting}[style=bashStyle]
docker build -t gohub:latest .
docker images gohub
# REPOSITORY  TAG     SIZE
# gohub       latest  8.2MB    <- Au lieu de 800+ Mo
\end{lstlisting}

% ============================================================
\moduleHeader{05}{L'Homme de l'Ombre --- Outils systeme avances (Bonus)}{GoGrey}
% ============================================================

\chapter{Reseau et Securite Systeme}

\section{Scanner de ports TCP concurrent}

\begin{lstlisting}[style=goStyle]
func Scan(ctx context.Context,
          hote string, debut, fin, concurrence int) []Resultat {

    resultats := make(chan Resultat, fin-debut+1)
    semaphore  := make(chan struct{}, concurrence) // Rate limiting
    var wg sync.WaitGroup

    for port := debut; port <= fin; port++ {
        wg.Add(1)
        p := port
        go func() {
            defer wg.Done()
            semaphore <- struct{}{}
            defer func() { <-semaphore }()

            portCtx, cancel :=
                context.WithTimeout(ctx, 300*time.Millisecond)
            defer cancel()
            resultats <- testerPort(portCtx, hote, p)
        }()
    }
    go func() { wg.Wait(); close(resultats) }()
    return collecterResultats(resultats)
}
\end{lstlisting}

\section{TLS natif et AES-GCM}

\begin{lstlisting}[style=goStyle]
// Serveur HTTPS sans OpenSSL
srv := &http.Server{
    TLSConfig: &tls.Config{MinVersion: tls.VersionTLS12},
}
srv.ListenAndServeTLS("server.crt", "server.key")

// Chiffrement AES-256-GCM (authentifie)
func Chiffrer(texte, cle []byte) ([]byte, error) {
    block, _ := aes.NewCipher(cle)  // 32 bytes = AES-256
    gcm, _   := cipher.NewGCM(block)

    nonce := make([]byte, gcm.NonceSize())
    io.ReadFull(rand.Reader, nonce)  // Nonce unique par appel !

    return gcm.Seal(nonce, nonce, texte, nil), nil
}
// Utilisez AES-GCM (authentifie) -- jamais AES-CBC
\end{lstlisting}

\chapter{Optimisation Bas Niveau}

\section{Le cycle d'optimisation}

\begin{center}
\begin{tikzpicture}[
  etape/.style={draw=GoBlue,fill=GoLight,rounded corners=4pt,
    minimum width=2.5cm,minimum height=0.7cm,font=\small\bfseries},
  arr/.style={->,>=Stealth,thick,color=GoBlue}]
\node[etape](m)  at (0,0)     {1. Mesurer};
\node[etape](id) at (3.5,0)   {2. Identifier};
\node[etape](c)  at (7,0)     {3. Comprendre};
\node[etape](o)  at (7,-1.8)  {4. Optimiser};
\node[etape](r)  at (3.5,-1.8){5. Re-mesurer};
\node[etape](rc) at (0,-1.8)  {6. Recommencer};
\draw[arr](m)--(id);\draw[arr](id)--(c);\draw[arr](c)--(o);
\draw[arr](o)--(r);\draw[arr](r)--(rc);
\draw[arr](rc) to[bend right=20] (m);
\end{tikzpicture}
\end{center}

\begin{boiteConseil}
\code{go test -bench=. -benchmem} mesure le temps ET les allocations.
La plupart des problemes de performance viennent des \textbf{allocations
memoire excessives}, pas du CPU. Mesurez avant d'optimiser.
\end{boiteConseil}

\begin{lstlisting}[style=goStyle]
// Pre-allouer : une seule allocation au lieu de N
donnees := make([]int, 0, n)

// sync.Pool : reutiliser des objets couteux
var pool = sync.Pool{
    New: func() any { return make([]byte, 0, 4096) },
}
buf := pool.Get().([]byte)
defer pool.Put(buf[:0])  // Remettre dans le pool propre
\end{lstlisting}

% ============================================================
%  ANNEXES
% ============================================================
\appendix

\chapter{Reference rapide}

\section{Outils de concurrence}

\begin{center}
\tableauStyle
\begin{tabularx}{\textwidth}{L >{\ttfamily\bfseries\color{GoBlue}}l}
\toprule
\rowcolor{GoDark}
\color{white}Besoin & \color{white}Outil \\
\midrule
Lancer une tache en arriere-plan   & go maFonction() \\
Passer des donnees entre goroutines & chan \\
Attendre plusieurs goroutines      & sync.WaitGroup \\
Proteger un etat partage           & sync.Mutex \\
Map thread-safe                    & sync.Map \\
Propager annulation / timeout      & context.Context \\
Attendre le premier channel pret   & select \\
Detecter les race conditions       & go run -race \\
\bottomrule
\end{tabularx}
\end{center}

\section{Toutes les commandes essentielles}

\begin{center}
\tableauStyle
\begin{tabularx}{\textwidth}{>{\ttfamily\small}l L}
\toprule
\rowcolor{GoDark}\color{white}Commande & \color{white}Description \\
\midrule
go run main.go              & Compiler et executer \\
go build -o app ./...       & Compiler en binaire \\
go build -ldflags="-s -w"   & Binaire optimise production \\
go test ./...               & Tous les tests \\
go test -race ./...         & Tests + race detector \\
go test -cover ./...        & Tests + couverture \\
go test -bench=. -benchmem  & Benchmarks avec memoire \\
go fmt ./...                & Formatter le code \\
go vet ./...                & Analyser les erreurs \\
go mod tidy                 & Nettoyer les dependances \\
CGO\_ENABLED=0 go build     & Binaire statique \\
GOOS=linux go build         & Cross-compilation Linux \\
\bottomrule
\end{tabularx}
\end{center}

% ============================================================
%  FIN
% ============================================================
\backmatter

\chapter*{A propos de l'auteur}
\addcontentsline{toc}{chapter}{A propos de l'auteur}

\begin{center}
\begin{tikzpicture}
  \fill[GoBlue](0,0) circle (1.2cm);
  \fill[GoDark](0,0) circle (0.78cm);
  \node[text=GoLight,font=\bfseries\Large] at (0,0) {AK};
\end{tikzpicture}

\vspace{5mm}
{\Large\bfseries\color{GoDark}Axel Konan}\\[2mm]
{\color{GoGrey}Architecte Logiciel
  $\cdot$ Formateur
  $\cdot$ Vulgarisateur Technique}\\[6mm]

\begin{tabular}{rl}
  \textbf{GitHub}    &
    \href{https://github.com/votre-pseudo}{github.com/votre-pseudo} \\
  \textbf{Formation} &
    \href{https://github.com/votre-pseudo/formation-go}%
         {github.com/votre-pseudo/formation-go} \\
  \textbf{Licence}   & Creative Commons CC BY-NC-SA 4.0 \\
\end{tabular}

\vspace{6mm}
\begin{minipage}{0.82\textwidth}
\centering\small\color{GoGrey}
Specialiste des systemes distribues et du Cloud Native,
Axel accompagne des equipes d'ingenieurs avec une approche
pragmatique~: partir du probleme reel avant d'introduire la theorie.
Son objectif~: rendre l'ingenierie de haut niveau accessible a tous.
\end{minipage}
\end{center}

\end{document}